% Paste into examplethesis.tex: replace from \chapter{Background} through the end of
% the first major section (keep later placeholder sections until 2.2+ are written).
% Cite keys map to Pre-Study References [1]--[6]. Comments note Pre-Study page anchors.

\chapter{Background}
\label{ch:background}

This chapter reviews research that underpins task-oriented assistance under vague, multi-turn input: how systems detect and resolve underspecification, how they represent uncertainty without relying only on surface text, how they model open-ended goals, how memory can personalise behaviour under data-minimisation constraints, and how such systems are evaluated.
The survey informs the modular design studied later in the thesis; detailed methods and experiments appear in subsequent chapters.

The remainder of the chapter is organised as follows.
First, active clarification under ambiguous instructions is reviewed.
Subsequent sections cover internal uncertainty handling and belief-oriented reasoning; open-ended goal inference and proactive user modelling; memory-augmented personalisation under privacy constraints; and evaluation paradigms for multi-turn intention tracking.
A closing summary synthesises remaining gaps and connects them to the research questions.

\section{From vague intent parsing to active clarification}
\label{sec:vague-clarification}
% Pre-Study §2.1 (p.~2); Refs.\ [1]--[6]

Traditional task-oriented pipelines often assume that users can state goals completely in a single utterance.
Empirical studies of human--LLM interaction challenge that assumption: even when natural-language interfaces feel effortless, users frequently produce prompts that are structurally incomplete or semantically ambiguous~\cite{koyuturk_understanding_2025}.
Under such conditions, treating each turn as a fully specified command is brittle; agents need mechanisms that recognise underspecification and decide whether to proceed, ask, or revise.

One response in the recent literature is to make clarification an explicit capability rather than an incidental side effect of generation.
Ask-to-Clarify, for example, introduces a detection stage that can interrupt execution when instructions are ambiguous and divert the dialogue into multi-turn clarification~\cite{lin_ask-to-clarify_2025}.
Beyond deciding \emph{whether} to ask, the content of clarification questions also matters.
Tang and Soulier~\cite{tang_clarifying_2025} show that prompting models to classify the type of ambiguity---for instance, distinguishing underspecification that calls for generalisation from cases that call for more precise constraints---before generating a query improves clarification relevance.
Together, these lines of work argue that ambiguity should be treated as a first-class control signal in dialogue, not merely as noise around an otherwise complete intent label.

Relying heavily on repeated clarification, however, carries documented costs.
Continual follow-up questions increase interaction effort and can undermine usability even when each individual question is locally informative.
Liu et~al.~\cite{liu_intent_2026} further report that large language models can degrade across turns when user constraints drift and intent alignment is lost---a failure mode they associate with multi-turn ``intent mismatch.''
In other words, clarification is not free: poorly timed questions may both annoy the user and destabilise the model's own dialogue state.

Cost-aware formulations therefore restrict when questions are asked.
Ghose et~al.~\cite{ghose_open-ended_2025} propose a bidirectional inference model that weighs expected information gain against user disruption, so that questions are issued only when the anticipated reduction in uncertainty outweighs interruption cost.
Fang et~al.~\cite{fang_uncertainty_2026} extend related ideas into a dual-mode joint planning setting, using dynamic programming to derive querying policies that trade off preference elicitation against interaction and verification costs.
The shared implication is that semantic disambiguation and user comfort must be optimised jointly: an agent that always clarifies may resolve ambiguity, yet still fail as a practical assistant.

Overall, the clarification literature establishes three linked requirements for systems that face vague booking-style requests: detect underspecification, generate type-appropriate questions when needed, and limit clarification by an explicit cost--benefit criterion.
The next section turns from outward questioning to inward uncertainty representation---how agents can quantify and act on residual ambiguity without always asking the user.
